\vspace{-3mm}
\section{Experimental Details.}
\label{app: experimental details}
\vspace{-3mm}

\noindent \textbf{Codebase}. Our implementation builds upon the original Duo~\href{https://github.com/s-sahoo/duo}{repository}~\citep{sahoo2025the} and the SEDD~\href{https://github.com/louaaron/Score-Entropy-Discrete-Diffusion}{repository}~\citep{lou2024discrete}, which serve as the primary codebases for our experiments. We extend these repositories to incorporate our proposed training framework, with the full algorithmic details provided in Appendix~\ref{app:algorithm}.

\noindent \textbf{Teacher checkpoints}. We use pretrained checkpoints for the distillation of MDLM and Duo models obtained from the Duo~\href{https://github.com/s-sahoo/duo}{repository}~\citep{sahoo2025the}, and a pretrained checkpoint for SEDD distillation from the SEDD~\href{https://github.com/louaaron/Score-Entropy-Discrete-Diffusion}{repository}~\citep{lou2024discrete}. Unfortunately, for SEDD, only a checkpoint corresponding to the absorbing process was available.

\noindent \textbf{Model architecture}. We follow~\citep{sahoo2025the} for distilling MDLM and Duo, and design the denoising model to operate on both continuous and discrete latents. We use a Transformer~\citep{vaswani2017attention} in which the first-layer token representations are computed via a matrix multiplication with the embedding matrix. For discrete inputs, we perform standard embedding lookups. In contrast, continuous inputs correspond to “soft lookups,” producing a convex combination of vocabulary embeddings. For SEDD distillation, we follow~\citep{lou2024discrete} and use a model that only accepts one-hot vectors. For all setups we use small models with $169$M parameters, specifically with $12$ layers, $12$ heads, hidden size $768$, sequence length $1024$, and dropout $0.1$.

\noindent \textbf{Training hyperparameters}.  
Across all experimental settings, we employed a per-device batch size of 8, resulting in a global batch size of 512. Optimization was performed using the AdamW optimizer~\citep{loshchilov2017decoupled}, with a fixed learning rate of $1 \times 10^{-6}$ for MDLM and Duo, and $3 \times 10^{-7}$ for SEDD. A constant learning rate schedule was used with 2500 warmup steps for all configurations. We adopted the log-linear noise formulation from~\citep{lou2024discrete} in the loss function, and applied exponential moving average (EMA) with a decay rate of 0.9999 on training.


\noindent \textbf{Evaluation protocol}. 
As observed by~\citep{zheng2024masked}, the GenPPL metric is sensitive to floating-point precision. To ensure consistency and numerical stability, we follow the protocol of~\citep{sahoo2025the} and perform all sampling experiments using \texttt{float64} precision. All evaluation metrics are computed using the official codebase provided in the Duo~\citep{sahoo2025the} repository.

\noindent \textbf{Dataset preparation}.
Following prior work~\citep{sahoo2024simple, sahoo2025the, lou2024discrete}, we preprocess the One Billion Words dataset using the detokenization procedure described by~\citet{lou2024discrete} and~\citet{sahoo2024simple}, with the official implementation available at \href{https://github.com/louaaron/Score-Entropy-Discrete-Diffusion/blob/main/data.py}{this link}. For the OpenWebText dataset, we utilize the \texttt{GPT2} tokenizer and apply a similar concatenation and wrapping procedure as in the previous works~\citep{sahoo2024simple, sahoo2025the}, targeting a sequence length of 1,024 tokens. During wrapping, \texttt{eos} tokens are inserted between consecutive sequences. As OpenWebText does not include an official validation split, we reserve the final 100{,}000 documents from the dataset as a validation set.

\noindent \textbf{Other baselines}.
Our primary baselines are SDTT~\citep{deschenaux2024beyond} and Duo-DCD~\citep{sahoo2025the}. Performance metrics for both methods are reported using values provided in the Duo~\citep{sahoo2025the}, corresponding to the best-performing results obtained after 5 rounds of distillation, described in~\citep{sahoo2025the}. While a valid pretrained checkpoint was available for Duo-DCD, we were unable to identify a usable checkpoint for SDTT~\citep{deschenaux2024beyond}, which precluded a fair comparison or distillation of this model within our experimental framework.